\section{Interpolation}

\subsection{Lagrange Interpolation}
\begin{formulabox}
$(n+1)$ points $\Rightarrow$ at most an $n$-th order polynomial:
\[
f(x) = \sum_{j=0}^{n} y_j \prod_{\substack{k=0 \\ k\neq j}}^{n} \frac{x - x_k}{x_j - x_k}
\]
\end{formulabox}

\subsection{Linear Interpolation}
\begin{formulabox}
Bracketing points $(x_2,y_2),(x_3,y_3)$, query $x_m$:
\[
\frac{y_3 - y_2}{x_3 - x_2} = \frac{y_m - y_2}{x_m - x_2}
\]
\[
\Longrightarrow\quad y_m = y_2 + \frac{y_3-y_2}{x_3-x_2}(x_m - x_2)
\]
Always use the two tabulated points that bracket the query $x$.
\end{formulabox}

\subsection{Newton's Divided-Difference Interpolation}
\begin{formulabox}
\[
P_n(x) = a_0 + a_1(x-x_0) + a_2(x-x_0)(x-x_1) + \dots
\]
\[
a_0 = y_0, \qquad a_1 = f[x_1,x_0] = \frac{y_1-y_0}{x_1-x_0}
\]
\[
a_2 = f[x_2,x_1,x_0], \quad \dots
\]
Master recurrence:
\[
f[x_0,x_1,\dots,x_{n+1}] = \frac{f[x_1,\dots,x_{n+1}] - f[x_0,\dots,x_n]}{x_{n+1}-x_0}
\]
Divided-difference table (top diagonal = coefficients $a_0,a_1,a_2,\dots$):
\[
\begin{array}{c c c c c}
x_0 & y_0 & & & \\
 & & f[x_0,x_1] & & \\
x_1 & y_1 & & f[x_0,x_1,x_2] & \\
 & & f[x_1,x_2] & & f[x_0,x_1,x_2,x_3]\\
x_2 & y_2 & & f[x_1,x_2,x_3] & \\
 & & f[x_2,x_3] & & \\
x_3 & y_3 & & &
\end{array}
\]
Advantage over Lagrange: adding a new point only appends one row/coefficient (no recomputation).
\end{formulabox}

\subsection{Spline Interpolation}
\begin{formulabox}
\begin{center}
\begin{tabular}{@{}lll@{}}
\toprule
Order & Piece & Unknowns \\
\midrule
Linear & $P_i = m_i x + c_i$ & $2n$ \\
Quadratic & $P_i = a_i x^2+b_ix+c_i$ & $3n$ \\
Cubic & $P_i = a_ix^3+b_ix^2+c_ix+d_i$ & $4n$ \\
\bottomrule
\end{tabular}
\end{center}
Quadratic equation count: $3n = \underbrace{2n}_{\text{through pts}} + \underbrace{(n-1)}_{\text{match }P'} + \underbrace{1}_{P_1''(x_1)=0}$

Cubic equation count: $4n = \underbrace{2n}_{\text{through pts}} + \underbrace{(n-1)}_{\text{match }P'} + \underbrace{(n-1)}_{\text{match }P''} + \underbrace{2}_{\text{boundary}}$

Natural cubic: $P_1''(x_{first})=0,\ P_n''(x_{last})=0$. Clamped cubic: end slopes $dy/dx$ given.
Runge's phenomenon $\Rightarrow$ motivation for splines (local, low-order pieces) over one global high-order polynomial.
\end{formulabox}
